\chapter{CƠ SỞ LÝ THUYẾT}

\section{Tổng quan bài toán phát hiện hành vi nguy hiểm}

\subsection{Định nghĩa và phạm vi bài toán}

Phát hiện hành vi nguy hiểm từ camera giám sát là bài toán thuộc lĩnh vực thị giác máy tính, nhằm tự động nhận dạng các tình huống đe dọa đến sức khỏe hoặc tính mạng con người thông qua phân tích luồng video theo thời gian thực. Trong phạm vi đề tài này, hành vi nguy hiểm được tập trung vào sự kiện té ngã — một trong những nguyên nhân hàng đầu gây thương tích nghiêm trọng ở người cao tuổi và bệnh nhân.

Bài toán được định nghĩa chính thức như sau: cho một chuỗi khung hình video $\{f_1, f_2, \ldots, f_T\}$ từ camera giám sát, hệ thống cần xác định tại mỗi thời điểm $t$ liệu người trong khung hình có đang thực hiện hành vi té ngã hay không, đồng thời phát cảnh báo tức thời khi phát hiện sự kiện.

\subsection{Các hướng tiếp cận hiện có}

Các phương pháp phát hiện té ngã hiện nay có thể phân thành ba nhóm chính:

\begin{itemize}
    \item \textbf{Dựa trên cảm biến (Sensor-based):} Sử dụng gia tốc kế, con quay hồi chuyển hoặc thiết bị đeo tay để đo chuyển động cơ thể. Ưu điểm là độ chính xác cao trong môi trường kiểm soát, nhưng nhược điểm là yêu cầu người dùng phải luôn mang thiết bị.

    \item \textbf{Dựa trên thị giác máy tính (Vision-based):} Phân tích hình ảnh từ camera mà không cần tiếp xúc trực tiếp. Hướng tiếp cận này đang được nghiên cứu rộng rãi nhờ sự phát triển của các kỹ thuật học sâu.

    \item \textbf{Kết hợp đa phương thức (Multimodal):} Tích hợp nhiều nguồn dữ liệu như camera, cảm biến âm thanh và cảm biến môi trường để tăng độ tin cậy.
\end{itemize}

Đề tài này tập trung vào hướng tiếp cận vision-based, kết hợp ước lượng tư thế người, phân tích luật và mô hình học sâu để đạt được sự cân bằng giữa độ chính xác và tốc độ xử lý thời gian thực.

\subsection{Thách thức và hướng giải quyết}

Bài toán phát hiện té ngã từ camera đối mặt với nhiều thách thức:

\begin{itemize}
    \item \textbf{Phân biệt té ngã với hành động tương tự:} Các hành động như ngồi xuống nhanh hoặc cúi người dễ bị nhầm lẫn với té ngã. Hệ thống kết hợp MediaPipe và YOLOv8-Pose để trích xuất keypoints cơ thể, tính vận tốc và góc nghiêng cơ thể theo thời gian thực. Mô hình Transformer được sử dụng để phân tích chuỗi chuyển động theo thời gian nhằm phân biệt quá trình té ngã với các hoạt động bình thường.


    \item \textbf{Biến đổi góc nhìn và chiếu sáng:} Camera đặt ở các góc độ khác nhau ảnh hưởng đến chất lượng ước lượng tư thế. Giải pháp là fusion hai mô hình MediaPipe và YOLOv8 để bù trừ lỗi.

    \item \textbf{Yêu cầu thời gian thực:} Hệ thống phải phản hồi trong vài giây để cảnh báo kịp thời. Giải pháp là tối ưu pipeline trên GPU và giảm độ trễ WebSocket.

    
\end{itemize}

% ─────────────────────────────────────────────────────────────────────────────

\section{Ước lượng tư thế người}

Ước lượng tư thế người (Human Pose Estimation) là kỹ thuật xác định vị trí các điểm khớp (keypoints) trên cơ thể người trong ảnh hoặc video, bao gồm vai, hông, gối, cổ tay và các khớp khác. Đây là nền tảng để phân tích chuyển động và nhận dạng hành vi.

\subsection{MediaPipe Pose}

MediaPipe Pose là framework ước lượng tư thế người được phát triển bởi Google, hoạt động theo thời gian thực với 33 điểm keypoint trên toàn thân. Framework sử dụng kiến trúc BlazePose — một mạng nơ-ron nhẹ được tối ưu cho thiết bị di động và máy tính cá nhân.

Quá trình xử lý của MediaPipe Pose gồm hai giai đoạn:

\begin{itemize}
    \item \textbf{Giai đoạn phát hiện (Detection):} Một mạng detector nhẹ xác định vùng chứa người trong khung hình.
    \item \textbf{Giai đoạn ước lượng (Estimation):} Mạng chính dự đoán tọa độ $(x, y, z)$ và độ tin cậy (visibility) của 33 keypoint.
\end{itemize}

Mỗi keypoint được biểu diễn bởi vector 4 chiều $(x, y, z, v)$ trong đó $x, y$ là tọa độ pixel đã chuẩn hóa, $z$ là độ sâu tương đối và $v \in [0, 1]$ là độ tin cậy. Trong đề tài này, MediaPipe được sử dụng như mô hình chính để trích xuất tư thế do khả năng xử lý nhanh và ổn định trên CPU.

\subsection{YOLOv8 Pose và ByteTrack}

YOLOv8 Pose là phiên bản mở rộng của kiến trúc YOLO (You Only Look Once) kết hợp phát hiện đối tượng và ước lượng tư thế trong một mô hình duy nhất. Mô hình dự đoán đồng thời bounding box của người và 17 keypoint theo chuẩn COCO.

Kiến trúc YOLOv8 gồm ba thành phần chính:
\begin{itemize}
    \item \textbf{Backbone:} Mạng trích xuất đặc trưng dựa trên CSPNet cải tiến.
    \item \textbf{Neck:} Mạng PAN-FPN tổng hợp đặc trưng đa tỉ lệ.
    \item \textbf{Head:} Đầu ra song song cho bounding box và keypoints.
\end{itemize}

ByteTrack là thuật toán theo dõi đa đối tượng (Multi-Object Tracking) được tích hợp vào YOLOv8 để duy trì ID ổn định cho mỗi người qua các khung hình. Thuật toán sử dụng IoU (Intersection over Union) kết hợp với bộ lọc Kalman để dự đoán vị trí và liên kết detection giữa các frame liên tiếp.

\subsection{Fusion đa mô hình}

Để tận dụng ưu điểm của cả hai mô hình, hệ thống thực hiện fusion có trọng số theo độ tin cậy:

\begin{equation}
    m_{fused} = \frac{w_{mp} \cdot m_{mp} + w_{yolo} \cdot m_{yolo}}{w_{mp} + w_{yolo}}
\end{equation}

trong đó $w_{mp}$ và $w_{yolo}$ lần lượt là độ tin cậy của MediaPipe và YOLOv8, $m$ là các chỉ số tư thế cần fusion (góc cơ thể, tỉ lệ bounding box). Vector đặc trưng thô được ghép nối từ 33 keypoint MediaPipe và 17 keypoint YOLO tạo thành vector 200 chiều đầu vào cho mô hình Transformer.

% ─────────────────────────────────────────────────────────────────────────────

\section{Phát hiện té ngã}

\subsection{Phương pháp dựa trên luật}

Phương pháp rule-based phát hiện té ngã dựa trên việc phân tích các đặc trưng hình học và vận tốc của cơ thể người theo thời gian. Phương pháp này có ưu điểm là có thể giải thích được (explainable) và không yêu cầu dữ liệu huấn luyện.

\subsection{Phân loại tư thế và phân tích vận tốc}

Hệ thống phân loại tư thế người thành bốn trạng thái: STANDING (đứng), SITTING (ngồi), LYING (nằm) và WALKING (đi lại). Việc phân loại dựa trên hai chỉ số chính:

\begin{itemize}
    \item \textbf{Góc cơ thể (body\_angle):} Góc giữa vector cột sống (từ vai đến hông) và trục thẳng đứng, tính theo công thức:
    \begin{equation}
        \theta = \arctan\left(\frac{|\Delta x_{spine}|}{|\Delta y_{spine}| + \epsilon}\right)
    \end{equation}
    Khi $\theta \geq 75°$, người được phân loại là LYING.

    \item \textbf{Tỉ lệ bounding box (aspect\_ratio):} Tỉ lệ chiều cao trên chiều rộng của bounding box $r = h/w$. Khi $r \leq 0.45$, người được phân loại là LYING.
\end{itemize}

Vận tốc hướng xuống của trọng tâm cơ thể được tính theo cửa sổ trượt:
\begin{equation}
    v_y = \frac{y_{center}(t) - y_{center}(t - \Delta t)}{\Delta t}
\end{equation}

Sự kiện té ngã được xác nhận khi thỏa mãn đồng thời các điều kiện: vận tốc đủ lớn ($v_y > v_{threshold}$), có lịch sử tư thế đứng/ngồi/đi, trạng thái LYING duy trì đủ số frame, thời gian chuyển trạng thái đủ ngắn, và không phải hành động cúi người.

\subsection{Mạng Transformer cho nhận dạng hành động}

Kiến trúc Transformer, được giới thiệu bởi Vaswani et al. (2017) trong bài báo ``Attention is All You Need'', ban đầu được thiết kế cho xử lý ngôn ngữ tự nhiên nhưng đã chứng minh hiệu quả vượt trội trong nhận dạng hành động từ chuỗi tư thế. Ưu điểm của Transformer so với RNN/LSTM là khả năng học mối quan hệ dài hạn giữa các frame thông qua cơ chế Self-Attention mà không bị giới hạn bởi độ dài chuỗi.

\textbf{Cơ chế Self-Attention} cho phép mô hình học được mối quan hệ giữa các frame trong chuỗi thời gian. Cho chuỗi đầu vào $X \in \mathbb{R}^{T \times d}$, Self-Attention được tính:

\begin{equation}
    \text{Attention}(Q, K, V) = \text{softmax}\left(\frac{QK^T}{\sqrt{d_k}}\right)V
\end{equation}

trong đó $Q = XW^Q$, $K = XW^K$, $V = XW^V$ là các ma trận Query, Key, Value; $d_k$ là chiều của không gian key. Multi-Head Attention kết hợp $h$ đầu attention song song để học nhiều mối quan hệ khác nhau:

\begin{equation}
    \text{MultiHead}(Q,K,V) = \text{Concat}(\text{head}_1, \ldots, \text{head}_h)W^O
\end{equation}

\textbf{Positional Embedding} là tham số học được $P \in \mathbb{R}^{(T+1) \times d_{model}}$ được cộng vào biểu diễn đầu vào để đưa thông tin vị trí thời gian vào mô hình, do Transformer không có tính chất tuần tự như RNN.

\textbf{CLS Token} là một vector đặc biệt $\mathbf{c} \in \mathbb{R}^{d_{model}}$ được thêm vào đầu chuỗi, học cách tổng hợp thông tin toàn chuỗi để phục vụ phân loại:

\begin{equation}
    \hat{y} = \text{MLP}(\text{TransformerEncoder}([\mathbf{c}; x_1; \ldots; x_T])[0])
\end{equation}

\subsection{Kết hợp Rule-based và AI}

Hệ thống sử dụng chiến lược AND logic — cả hai thành phần phải đồng ý mới phát cảnh báo:

\begin{equation}
    \text{fall\_detected} = \text{rule\_fall} \wedge \text{ai\_fall}
\end{equation}

Trong đó $\text{rule\_fall}$ là kết quả từ FallDetector (rule-based) và $\text{ai\_fall}$ là kết quả từ TransformerEngine với ngưỡng xác suất $p_{fall} \geq 0.65$. Chiến lược này giúp giảm cả false positive lẫn false negative so với dùng từng phương pháp riêng lẻ.

\subsection{Đánh giá mô hình với K-Fold Cross Validation}

K-Fold Cross Validation là kỹ thuật đánh giá mô hình học máy nhằm giảm thiểu sự phụ thuộc vào cách phân chia tập dữ liệu, đặc biệt hữu ích khi dữ liệu có kích thước giới hạn.

Quy trình K-Fold Cross Validation gồm các bước sau:

\begin{enumerate}
    \item Chia toàn bộ dataset $\mathcal{D}$ thành $K$ phần (fold) có kích thước xấp xỉ bằng nhau: $\mathcal{D} = \{D_1, D_2, \ldots, D_K\}$.
    \item Lặp $K$ lần: ở vòng lặp thứ $i$, sử dụng $D_i$ làm tập validation và $K-1$ phần còn lại làm tập huấn luyện.
    \item Huấn luyện và đánh giá mô hình tại mỗi fold, thu được độ chính xác $\text{acc}_i$.
    \item Tính độ chính xác trung bình và độ lệch chuẩn:
    \begin{equation}
        \overline{\text{acc}} = \frac{1}{K}\sum_{i=1}^{K}\text{acc}_i, \qquad
        \sigma = \sqrt{\frac{1}{K}\sum_{i=1}^{K}(\text{acc}_i - \overline{\text{acc}})^2}
    \end{equation}
\end{enumerate}

Ưu điểm của K-Fold so với phân chia train/test đơn giản:

\begin{itemize}
    \item \textbf{Đánh giá tin cậy hơn:} Mỗi mẫu dữ liệu được dùng làm validation đúng một lần, tránh sự ngẫu nhiên của một lần phân chia.
    \item \textbf{Phát hiện overfitting:} Độ lệch chuẩn $\sigma$ cao giữa các fold cho thấy mô hình không ổn định và có thể đang overfit trên một phần dữ liệu.
    \item \textbf{Tận dụng toàn bộ dữ liệu:} Đặc biệt quan trọng khi tập dữ liệu nhỏ, vì tất cả các mẫu đều được dùng cho cả huấn luyện lẫn đánh giá.
\end{itemize}

Đề tài sử dụng Stratified K-Fold với $K = 5$, đảm bảo tỉ lệ nhãn \textit{fall}/\textit{non\_fall} được giữ nguyên trong mỗi fold để tránh mất cân bằng dữ liệu ảnh hưởng đến kết quả đánh giá. Ngoài ra, kỹ thuật data augmentation (lật ngang, thay đổi tốc độ, thêm nhiễu Gaussian) chỉ được áp dụng trên tập huấn luyện của từng fold, không áp dụng trên tập validation, đảm bảo đánh giá khách quan.

% ─────────────────────────────────────────────────────────────────────────────

\section{Nhận diện khuôn mặt}

\subsection{RetinaFace và ArcFace}

RetinaFace là mô hình phát hiện khuôn mặt (face detection) đa tác vụ, đồng thời dự đoán bounding box khuôn mặt, 5 điểm landmark (mắt, mũi, miệng) và điểm tin cậy. Mô hình sử dụng kiến trúc Feature Pyramid Network để phát hiện khuôn mặt ở nhiều tỉ lệ khác nhau, phù hợp với môi trường camera giám sát có nhiều người ở các khoảng cách khác nhau.

ArcFace (Additive Angular Margin Loss) là phương pháp học biểu diễn khuôn mặt (face representation) hiện đại, huấn luyện mạng để tạo ra vector đặc trưng 512 chiều cho mỗi khuôn mặt. ArcFace thêm margin góc $m$ vào hàm loss để tăng khả năng phân biệt giữa các danh tính:

\begin{equation}
    \mathcal{L} = -\frac{1}{N}\sum_{i=1}^{N}\log\frac{e^{s\cos(\theta_{y_i} + m)}}{e^{s\cos(\theta_{y_i}+m)} + \sum_{j \neq y_i} e^{s\cos\theta_j}}
\end{equation}

trong đó $\theta_{y_i}$ là góc giữa vector đặc trưng và trọng số lớp, $s$ là hệ số tỉ lệ và $m$ là margin góc.

\subsection{Cosine Similarity}

Để so sánh hai vector đặc trưng khuôn mặt $\mathbf{u}$ và $\mathbf{v}$, hệ thống sử dụng độ tương đồng cosine:

\begin{equation}
    \text{sim}(\mathbf{u}, \mathbf{v}) = \frac{\mathbf{u} \cdot \mathbf{v}}{\|\mathbf{u}\| \cdot \|\mathbf{v}\|}
\end{equation}

Giá trị $\text{sim} \in [-1, 1]$, trong đó giá trị càng gần 1 thì hai khuôn mặt càng giống nhau. Ngưỡng nhận dạng được đặt tại $\text{sim} \geq 0.6$ để cân bằng giữa false accept rate và false reject rate.

% ─────────────────────────────────────────────────────────────────────────────

% ─────────────────────────────────────────────────────────────────────────────

\section{Các công nghệ sử dụng}

\subsection{FastAPI}

FastAPI là framework Python hiện đại để xây dựng REST API, dựa trên tiêu chuẩn OpenAPI và JSON Schema. FastAPI sử dụng type hints của Python và Pydantic để tự động xác thực dữ liệu đầu vào, đồng thời tự động sinh tài liệu API tương tác (Swagger UI). Framework hỗ trợ lập trình bất đồng bộ (async/await) với ASGI, cho phép xử lý hàng nghìn kết nối đồng thời với hiệu năng cao.

\subsection{Flutter}

Flutter là framework phát triển ứng dụng đa nền tảng của Google, sử dụng ngôn ngữ Dart. Flutter biên dịch trực tiếp sang mã máy native, cho hiệu năng gần tương đương ứng dụng native. Hệ thống widget phong phú cho phép xây dựng giao diện đẹp, nhất quán trên cả Android và iOS từ một codebase duy nhất. Ứng dụng mobile trong đề tài sử dụng Provider cho quản lý trạng thái và GoRouter cho điều hướng.

\subsection{NuxtJS}

NuxtJS là framework Vue.js cho phát triển ứng dụng web, hỗ trợ Server-Side Rendering (SSR) và Static Site Generation (SSG). Trong đề tài, NuxtJS được dùng để xây dựng giao diện web quản trị cho phép admin theo dõi và quản lý hệ thống từ trình duyệt.

\subsection{PostgreSQL và Neon}

PostgreSQL là hệ quản trị cơ sở dữ liệu quan hệ mã nguồn mở, nổi tiếng với tính ổn định, hỗ trợ ACID và khả năng mở rộng. Neon là dịch vụ PostgreSQL serverless trên đám mây, cho phép tự động mở rộng theo tải và tách biệt môi trường đọc/ghi để tối ưu hiệu năng. Hệ thống sử dụng SQLAlchemy với chế độ async (asyncpg) để truy vấn không đồng bộ.

\subsection{Firebase Authentication}

Firebase Authentication là dịch vụ xác thực của Google, cung cấp cơ chế đăng nhập an toàn với nhiều phương thức như email/mật khẩu và OAuth. Hệ thống sử dụng JWT (JSON Web Token) từ Firebase để xác thực người dùng tại backend, đảm bảo mỗi request đều được kiểm tra quyền truy cập.

\subsection{Cloudinary}

Cloudinary là dịch vụ lưu trữ và quản lý media trên đám mây, hỗ trợ upload, chuyển đổi và phân phối ảnh/video qua CDN toàn cầu. Khi phát hiện té ngã, hệ thống cắt đoạn video 6 giây (90 frame trước + 90 frame sau sự kiện) và upload lên Cloudinary để lưu trữ và chia sẻ qua URL với người thân.

% \subsection{ONNX Runtime}

% ONNX (Open Neural Network Exchange) là định dạng chuẩn mở để lưu trữ và triển khai mô hình học máy. ONNX Runtime là engine tối ưu hóa thực thi mô hình ONNX, hỗ trợ tăng tốc trên nhiều loại phần cứng (CPU, GPU, NPU). Trong hệ thống, ONNX Runtime được sử dụng để triển khai mô hình nhận diện khuôn mặt (RetinaFace/ArcFace) với hiệu năng cao mà không cần framework học sâu đầy đủ.